% !PRINTED 82
\weberpage{82}
%<<CONT

\begin{equation}
\begin{aligned}
 a_1^{(1)}x_1+a_2^{(1)}x_2+\cdots+a_n^{(1)}x_n&=0,\\
 a_1^{(2)}x_1+a_2^{(2)}x_2+\cdots+a_n^{(2)}x_n&=0,\\
 &\ \vdots\\
 a_1^{(m)}x_1+a_2^{(m)}x_2+\cdots+a_n^{(m)}x_n&=0.
\end{aligned}
\tag{1}
\end{equation}
worin die Coëfficienten $a_i^{(k)}$ als gegebene Grössen betrachtet werden. Ueber die Zahlen $m,n$ wollen wir vorläufig noch gar keine Voraussetzung machen, sondern uns allgemein die Aufgabe stellen, \gesperrt{alle Werthsysteme} der $x_1,x_2,\ldots,x_n$ zu ermitteln, die den Gleichungen (1) genügen.

Eine Lösung der Gleichungen (1) können wir sofort angeben; sie sind nämlich, was auch die Coëfficienten $a_i^{(k)}$ sein mögen, erfüllt, wenn
\begin{equation}
 x_1=0,\quad x_2=0,\quad \ldots,\quad x_n=0.
\tag{2}
\end{equation}
Einen anderen extremen Fall können wir noch erwähnen: wenn nämlich die Coëfficienten $a_i^{(k)}$ sämmtlich den Werth Null haben, dann sind die Gleichungen (1) für beliebige Werthe von $x_1,x_2,\ldots,x_n$ befriedigt.

Der allgemeinen Beantwortung der Frage schicken wir folgende Bemerkungen voraus. Wir schreiben das System der Coëfficienten von (1) in Form eines Rechtecks
\begin{equation}
\begin{array}{cccc}
 a_1^{(1)}&a_2^{(1)}&\cdots&a_n^{(1)}\\
 a_1^{(2)}&a_2^{(2)}&\cdots&a_n^{(2)}\\
 \vdots&\vdots&&\vdots\\
 a_1^{(m)}&a_2^{(m)}&\cdots&a_n^{(m)}.
\end{array}
\tag{3}
\end{equation}
Ein solches Schema, das für sich noch keine numerische Bedeutung hat, heisst eine \gesperrt{Matrix}, insofern es als Quelle einer grösseren Anzahl von Determinanten betrachtet wird. \gesperrt{Die der Matrix entstammenden Determinanten erhält man, wenn man beliebige Zeilen und Colonnen weglässt, in beliebiger, nur insoweit bestimmter Anzahl, dass die übrig bleibenden Elemente ein Quadrat bilden, und dieses Quadrat als Determinante auffasst.} So erhält man aus der Matrix einreihige, zweireihige u. s. f. Determinanten. Die höchsten Determinanten sind $n$- oder $m$-reihig, je nachdem $n$ oder $m$ die kleinere Zahl ist (oder $n$-reihig, wenn $n=m$ ist).
